\chapter{正交集方法在矩阵计算上的应用}

\section{本章简介}

矩阵代数中的常用技术根植于线性等式的概念,并且仅在其\underline{原始}意义上对矩阵的元素,行和列进行操作.关于使用标准方法求解矩阵逆矩阵问题的完整论述,本章不做过多阐述.本章发展的技术基于\underline{对偶},\underline{极性}或\underline{正交}概念,这些概念与线性不等式相关,属于极性理论和择一性理论.正如我们将看到的,对偶(正交)方法除了非常有用之外,还在不增加复杂性的情况下拓展了线性等式的应用领域.

在本章中,我们将展示如何应用第1章中描述的正交集来解决矩阵问题.第2节提出了一种功能强大的矩阵求逆方法.正交分解的视角有助于理解整个正交化过程,并且正如第3节和第4节所分别展示的,它还可以用于计算初始矩阵的秩和行列式.在第5节中,我们给出了一个用于获取矩阵的秩,计算行列式以及在矩阵可逆时求其逆的算法.第6节中我们计算了该方法的复杂度,结果表明它与高斯消元法的复杂度一致.在第7节,我们展示了当原始矩阵的少数几行和(或)几列发生改变时,如何以极少的额外计算开销来更新逆矩阵和行列式.

大多数现有的计算矩阵逆或行列式的方法,在原始矩阵的单个元素或几个元素发生改变时,都需要从头开始重新执行整个过程.因此,如果对原始矩阵做出修改,之前这些计算所付出的大量计算努力就白费了.然而,还有一些替代方法,例如Sherman-Morrison方法和Woodbury方法,它们允许在对初始矩阵进行某些特殊类型的修改时,以减少乘法和加法次数的方式来计算新的逆矩阵.这些方法起源于矩阵理论,而本文提出的方法可以直接应用于这种动态更新过程.

在第8节中,我们表明了所提出的方法适用于计算符号逆矩阵或符号行列式,即处理以符号或参数作为元素(而非具体数值)的矩阵.事实上,结果表明,这个问题可以通过数值程序轻松解决,而这种求逆方法极大地简化了数值程序.在第9节中,我们把这一方法推广到分块矩阵.在第10节中,我们给出了在增加和(或)删除行和(或)列后更新矩阵逆的算法.

总结来说,在本章中,我们提供了一种求逆矩阵,计算行列式和获取矩阵秩的方法,这种方法计算效率高,并允许在初始矩阵的行发生改变时对逆和行列式进行动态更新.此外,我们表明该方法特别适用于符号计算.这种方法是针对矩阵的列操作定义的,但它同样适用于行操作.最后,我们指出了该方法适用于分块定义的矩阵,并具有易于且显而易见的并行实现的额外优势.